QUIC is a general-purpose transport protocol, first standardized by the Internet Engineering Task Force (IETF) in 2021, which is developed on top of UDP and combines features from mature protocols such as TCP and TLS. QUIC differentiates itself in that it offers secure connections by default, promises fewer round trips before data is transmitted, fixes the head-of-line blocking present with TCP, and is designed to be extensible through its user space implementations. However, previous work has shown packet I/O to be QUIC's largest bottleneck in terms of its maximum achievable throughput in high-bandwidth networks, and many QUIC implementations fail to reach the same goodputs as TCP+TLS in such environments. This highlights a need for further optimization of the UDP+QUIC datapath to match TCP's performance.

In this thesis, we investigate adding \emph{EXpress Data Path (XDP)} to the QUIC datapath to analyze its impact on latency and reducing packet I/O. We use the MsQuic implementation from Microsoft as it offers support for XDP. We show that while XDP is not able to acheive a higher throughput than MsQuic with the UDP socket optimizations \emph{message batching} and \emph{segmentation offloading}, it is significantly higher than the MsQuic datapath with no optimizations. To analyze latency, we sync the hardware clocks and system time between a packet sniffer and client using Precision Time Protocol (PTP) and analyze the time difference between receiving the packet in userspace within the client and receiving the packet on the NIC in the sniffer. We show that there is no significant difference between the XDP and non-XDP cases. Finally, we create a taxonomy in which we evaluate various QUIC optimizations addressing the packet I/O limitation, organizing each based on implementation complexity, performance improvement potential, unwanted side effects, flexibility, privilege level usage, and deployment.

% sentence 1 introduces the topic of the paper so it's clear what QUIC builds off of.
% sentence 2 summarizes work 
% sentence 3 

% QUIC differentiates itself in that it offers secure connections by default, promises fewer round trips before data is transmitted, fixes the head-of-line blocking present with TCP, and is designed to be extensible through its user space implementations. These features have made QUIC an ideal foundation for new application-layer protocols such as Hypertext Transfer Protocol 3 (HTTP/3)~\cite{rfc9114}, which now commands 20.5\% of HTTP traffic worldwide~\cite{cloudflare2024}. This underscores the importance of optimizing QUIC’s performance as its adoption and significance continue to grow. Relatedly, previous work from Jaeger et al.~\cite{quic_on_the_highway} has shown that packet I/O is QUIC's largest bottleneck in terms of its maximum achievable throughput. The following paragraphs explore a few approaches to addressing this limitation.

% EXpress Data Path (XDP) is a high performance data path in the Linux kernel based on the extended Berkeley Packet Filter (eBPF) runtime. First introduced in 2016~\cite{linux_xdp_merge_2016}, XDP works by attaching a hook into the network receive (RX) path, which can be used to handle packets at the earliest possible time, before any other work is done. This addition is a native Linux feature and adds an entire class of functionality that simply did not exist beforehand. With AF\_XDP, a socket type, the range of possibilities is extended yet again to allow sending raw layer 2 frames directly to the user space QUIC implementations. At the cost of some additional header processing, AF\_XDP allows the packets to skip the Linux network stack entirely, which adds much overhead to each packet. Moreover, a Linux kernel patch from 2023 enables Ethernet frame checksum offloading specific to AF\_XDP~\cite{fomichev2023xsk_tx_metadata}, which shifts its computation from the CPU to the Network Interface Controller (NIC), corresponding in less CPU usage.

% Another approach to the packet I/O limitation is DPDK~\cite{dpdk}, a software framework providing a set of data plane libraries and NIC polling-mode drivers to offload all packet processing from kernel space into user space. The fundamental difference between XDP and DPDK is that XDP represents network stack bypass while DPDK represents kernel bypass. This allows DPDK to avoid overheads generated by context switching, copying data between kernel and user space, interrupt handling, and the Linux scheduler. However, this is accomplished at the expense of having to reimplement many of the Linux network stack features such as firewalling, traffic shaping, routing, and socket APIs.

% In-kernel QUIC~\cite{quic_lxin} is yet another approach, which moves QUIC from user space into kernel space. Advantages include seamless integration with kernel subsystems, standardized socket APIs for QUIC, and reduced context switches. However, this also comes at the cost of slow update cycles due to the nature of conservative Linux kernel upgrades, which user space implementations do not suffer from.

% In sum, this Master's Thesis shall investigate the performance improvements in throughput and reduced CPU usage as a result of combining AF\_XDP with QUIC and patching MsQuic to add support for Ethernet frame checksum offloading. This is done in parallel to creating a taxonomy in which we evaluate various QUIC optimizations addressing the packet I/O limitation, organizing each based on implementation complexity, performance improvement potential, unwanted side effects, flexibility, privilege level usage, and dynamic vs. static loading.

% The first sentence of an abstract should clearly introduce the topic of the paper so that readers can relate it to other work they are familiar with.
% However, an analysis of abstracts across a range of fields show that few follow this advice, nor do they take the opportunity to summarize previous work in their second sentence.
% A central issue is the lack of structure in standard advice on abstract writing, so most authors don’t realize the third sentence should point out the deficiencies of this existing research.
% To solve this problem, we describe a technique that structures the entire abstract around a set of six sentences, each of which has a specific role, so that by the end of the first four sentences you have introduced the idea fully.
% This structure then allows you to use the fifth sentence to elaborate a little on the research, explain how it works, and talk about the various ways that you have applied it, for example to teach generations of new graduate students how to write clearly.
% This technique is helpful because it clarifies your thinking and leads to a final sentence that summarizes why your research matters.

% \textit{The text above describes the contents and structure of an abstract.
% The text by Steve Easterbrook~\cite{abstractwriting} presents the structure of typical abstract for a scientific paper.
% The abstracts of theses are typically not as space-constraint as scientific papers, therefore, abstracts tend to get longer.}
